\section{Critérios de design e comunicação}
\label{sec:design}

O projeto visual combina princípios de aprendizagem multimídia e de comunicação persuasiva. A teoria cognitiva da aprendizagem multimídia sustenta que as pessoas aprendem melhor com palavras e imagens em conjunto do que com palavras isoladas, e que a atenção deve ser protegida de elementos supérfluos \cite{mayer2009}. Em consequência, adotaram-se os seguintes critérios:

\begin{alineas}
  \item uma mensagem principal por lâmina, enunciada no título;
  \item texto mínimo: frases curtas, números em destaque e ícones no lugar de parágrafos, com a explicação oral reservada ao roteiro das notas do apresentador;
  \item hierarquia tipográfica clara, com títulos grandes e corpo legível a distância, e alto contraste entre texto e fundo;
  \item paleta restrita, com uma cor dominante (verde-petróleo), uma cor de destaque (coral) e neutros, aplicada em todas as lâminas;
  \item gráficos nativos do \textit{software} de apresentação, editáveis pelo grupo, com valores rotulados;
  \item toda cifra com fonte, e toda projeção rotulada como estimativa ou hipótese;
  \item identificação nominal dos integrantes na primeira e na última lâmina, conforme exigido pela orientação.
\end{alineas}

A apresentação não utiliza fotografias de pessoas reais ou de animais de terceiros, de modo a evitar questões de direito de imagem e de consentimento; o apelo visual vem de ícones, formas e gráficos. Caso o grupo disponha de fotografias próprias, com autorização, podem ser acrescentadas às lâminas 2 e 4 sem alterar a estrutura.
